\section{Results and Conclusions}
To conclude, it is worth highlighting on the successful creation of a scalable and modular benchmark for evaluating the PDDL planning capabilities of Large Language Models. The framework was designed to be extensible to a wide range of problems and models, with a strong emphasis on producing easily interpretable and traceable results. The \textit{raw, parsed, scored} output structure ensures diagnostic precision, allowing for a clear separation between generation, parsing, and validation failures.

\subsection{General Conclusions on LLM Planning Capabilities}
Our evaluation of several models highlighted critical insights about the current state of LLM-based planning: the main finding is that while LLMs exhibits a complete understanding of PDDL syntax, their ability to perform robust, long-horizon, and goal-directed reasoning is still fragile.
Models frequently generate action sequences that are syntactically correct and locally executable, as pointed out by generally high \texttt{Executability Ratios} and low \texttt{Hallucination Rates}, but these plans often fail to achieve the global goal, leading to low overall \texttt{Success Rates}. The significant gap between first-attempt success (\texttt{FASR}) and the final success rate (\texttt{SR}) highlights a strong dependency on iterative repair; this suggests that current models engage in a behavior more prone to stochastic search rather than more deterministic ones.

\subsection{Limitations of the Study}
It is worth noticing also limitations of this study: the evaluation could be still contrained to the available hardware, since these LLMs models require significant resources to run effectively.
Also the observed performance is inherently tied to the specific prompting written by us and iterative repair strategies employed; different approaches could yield different results.

\subsection{Hints for future works}
The modularity of the benchmark creates everal possibilities for future work and improvement:
\begin{itemize}
    \item \textbf{Single Domain Testing:} The set of PDDL problems could be reduced to a single Reasoning type and the difficulty increased graduelly to see where a given LLM starts to fail the planning.
    
    \item \textbf{Advanced Interaction Protocols:} Future work could explore more sophisticated interaction protocols. This includes implementing more informative feedback mechanisms for the iterative repair loop, providing models with specific details about precondition failures (e.g., "action `(A)` failed because precondition `(P)` was false") to better assess their self-correction abilities.
    
    \item \textbf{Broader Model Evaluation:} The benchmark should be used to evaluate a broader and more diverse set of models, to build a more comprehensive landscape of LLM planning skills.

    99\item \textbf{Single Model with different releases Evaluation:} The benchmark could be used to evaluate a single models but different versions of it, to test how good the capabilites have grown each new release of the LLM.
    
    \item \textbf{Deeper Analytical Insights:} The analytical layer can be extended to perform deeper correlational studies like linking specific model capability profiles to performance on domains with certain structural properties (e.g., high branching factor or sparse rewards) could reveal fundamental strengths and weaknesses of different architectures.
\end{itemize}
